\section{Compliance and national policy incentives}\label{sec:policy}

We begin with the continuation game after private investment. Two distinctions are central. A firm that has already invested in the open family faces a different incremental cost from a firm that has not invested. Moreover, a government's desire to replace a particular private standard is not sufficient to infer that such a standard can arise in a pure equilibrium of the \emph{preceding} investment game.

\subsection{A public mandate is implemented}\label{subsec:adoption}

Following a public mandate, an active compliant firm uses \(O(1,0)\); a noncompliant firm exits. With no other compliant producers, the firm's gross Cournot return is \(3/(4m)\). When one or two other firms comply, its operating return is, respectively,
\begin{equation}
 \pi_I^{(2)}
 =3m\left(\frac{1+\rho}{2m+\gamma_0-v}\right)^2,
 \quad
 \pi_I^{(3)}
 =3m\left(\frac{1+\rho}{2(m+\gamma_0-v)}\right)^2.
 \label{eq:adoption-profits}
\end{equation}
On (\ref{eq:baseline-box}), the isolated-adopter return is the smallest of the three. No initial history can require an \(N\) firm to spend more than \(K+\phi/2\) to comply, or a previous \(O\) firm to spend more than \(\phi/2\) in additional outlays. Since
\begin{equation}
 \frac{3}{4m}=\frac{25}{33}
 >\frac{14}{25}+\frac{7}{100},
 \label{eq:adoption-margin}
\end{equation}
adoption is strictly preferable to exit regardless of the other firms' actions.

\begin{lemma}[Mandatory compliance]\label{lem:adoption}
For every initial technology profile and every parameter in
(\ref{eq:baseline-box}), all firms' compliance with an enacted public
mandate is the unique Nash continuation of the compliance subgame.
\end{lemma}
\begin{proof}
Equation (\ref{eq:adoption-profits}) computes the operating payoff for
each possible number of other adopters. Inequality
(\ref{eq:adoption-margin}) compares their minimum with the largest
incremental installation cost faced by any firm. Sunk costs do not
change the payoff from compliance versus exit. Hence compliance is
strictly dominant for each firm, and its unique Nash profile consists
of three adopters.
\end{proof}

The induced public Cournot quantities and national welfare are therefore those in (\ref{eq:qi})--(\ref{eq:WI}), independently of the prior private designs. Let \(\mathcal W_i^L(h)\) denote country \(i\)'s payoff if intervention does not occur at the observed initial history \(h\). Define its net policy gain
\begin{equation}
 \Delta_i(h)=\mathcal W_i^M(h)-\mathcal W_i^L(h).
 \label{eq:policy-gain}
\end{equation}
Our explicit motion and tie conventions imply the public mandate passes if and only if
\begin{equation}
 \left|\{i:\Delta_i(h)\ge0\}\right|\ge r
 \quad\text{and}\quad
 \calA\cap\{i:\Delta_i(h)\ge0\}\ne\varnothing.
 \label{eq:vote-rule}
\end{equation}
This simple expression will help isolate the role of agenda power. It applies \emph{at a fixed history}; later we must ask which histories are privately supportable.

\subsection{Histories with at most one private developer}\label{subsec:zero-one}

When the initial profile has no developers or only one, there is no compatible private pair. The Cournot matrix remains at its all-\(N\) values, even if one firm has sunk an open-technology development cost. Let \(W_N^g\) be one country's initial gross firm-profit-plus-consumer-surplus value when all products have \(z_i=1\) and off-diagonal \(M_{ij}=\gamma_0\). Then
\begin{equation}
 q_N=\frac{1}{2(m+\gamma_0)},\quad
 W_N^g=3m q_N^2+
 \frac{3}{2}(m+2\gamma_0)q_N^2.
 \label{eq:WN}
\end{equation}
An initially \(N\) country's net policy gain from \(M\) is
\[
 \Delta_N=W_I-K-W_N^g.
\]
For an initial \(O(c,a)\) developer, comparing the two \emph{same-sunk-cost} histories yields
\[
 \Delta_O=W_I-W_N^g+\frac{\phi c^2}{2}
         =\Delta_N+K+\frac{\phi c^2}{2}.
\]
Both gains are strictly positive throughout (\ref{eq:baseline-box}), since \(W_I-W_N^g>K_{\max}=.56\). Thus all governments vote YES after zero or one private developer, irrespective of the constitution's nonempty agenda set and quorum.

\begin{lemma}[Unanimous demand for intervention at sparse histories]
\label{lem:sparse}
After any initial profile with zero or one privately developed open
interface, the three national governments strictly prefer the public
mandate. Hence every nonempty agenda set and every quorum \(r\le3\)
enacts \(M\).
\end{lemma}
The strict comparison is the same after any feasible \(O(c,a)\)
singleton investment: private compatibility externalities are absent
until a second open developer arrives.

\subsection{A privately optimal bilateral regional standard}\label{subsec:bilateral}

Consider firms 1 and 2 developing a common \(O(1,a)\), while firm 3 remains undeveloped. Set
\begin{equation}
 b=\gamma_0+\kappa a-v,\qquad d=\gamma_0-\eta a,
 \qquad z=1+\rho,\qquad
 D=2m(2m+b)-2d^2.
 \label{eq:bilateral-definitions}
\end{equation}
Solving (\ref{eq:cournot-foc}) for the two insiders and the outsider gives
\begin{equation}
 x(a,\eta)=\frac{2mz-d}{D},
 \qquad
 y(a,\eta)=\frac{2m+b-2dz}{D}.
 \label{eq:xy}
\end{equation}
Here \(x\) is either founder's quantity in each national market, and \(y\) is the outsider's. National consumer surplus is
\begin{equation}
 C_R(a,\eta)
 =(m+b)x^2+2dxy+\frac{m}{2}y^2.
 \label{eq:CR}
\end{equation}
The founder's total private return before the fixed setup \(K\), and the outsider's gross profit, are
\begin{equation}
 f(a,\eta)=3m x^2-\frac{\phi}{2}(1+a^2),
 \qquad
 \pi_O(a,\eta)=3m y^2.
 \label{eq:private-returns}
\end{equation}

For our parameter rectangle, an installed bilateral open cohort chooses full compatibility \(c=1\) and a unique interior commonality level
\begin{equation}
 a_R^*(\eta)\in(.25,.61),\quad
 f_a(a_R^*(\eta),\eta)=0,\quad
 f_{aa}<0,\quad \frac{da_R^*}{d\eta}>0.
 \label{eq:a-star}
\end{equation}
These claims concern the constrained bilateral design maximization, not the complete six-player game. The global one-dimensional optimum, including all design-boundary checks, is verified in Appendix~\ref{app:bilateral}.

At this observed optimal regional standard, the national policy gains take the form
\begin{equation}
 (\Delta_1,\Delta_2,\Delta_3)
 =\left(B(\eta),B(\eta),K_V(\eta)-K\right),
 \label{eq:bilateral-votes}
\end{equation}
where
\begin{align}
 B(\eta)
 &=W_I-\frac{\phi}{2}a_R^*(\eta)^2
       -f(a_R^*,\eta)-C_R(a_R^*,\eta),
       \label{eq:B}\\
 K_V(\eta)
 &=W_I-\pi_O(a_R^*,\eta)-C_R(a_R^*,\eta).
       \label{eq:KV}
\end{align}
The two founders' governments strictly prefer the public mandate, \(B(\eta)>0\), while the excluded government's gain \(K_V(\eta)-K\) is strictly negative on (\ref{eq:baseline-box}). The underlying inequalities are certified uniformly over the entire privately optimal design enclosure, not inferred from a single numerical illustration. They deliver exactly two favorable governments. Whether that majority can translate its preferences into policy depends on the proposal entitlement set.
